\section{Gêmeos Digitais na Era dos Ecossistemas Abertos}

Um dos desenvolvimentos paradigmáticos mais cruciais documentados pela literatura científica contemporânea é a convergência simbiótica entre gêmeos digitais e \textbf{ecossistemas abertos de IA}. Historicamente, a adoção de tecnologias de simulação médica foi freada por altos custos de licenciamento e pelas barreiras intransponíveis dos ecossistemas proprietários (\textit{vendor lock-in}). 

Iniciativas disruptivas baseadas em \textit{open-source}, como o \textbf{OpenCode Ecosystem}~\cite{opencode2026ecosystem} e frameworks composicionais como o \textbf{OpenTwins}~\cite{robles2023opentwins, infante2024integrating}, emergiram como a solução definitiva para democratizar o acesso à saúde digital de vanguarda. Estas plataformas operam sob contratos de interoperabilidade universais, permitindo que pesquisadores e clínicos orquestrem simulações biomecânicas avançadas sem a necessidade de infraestruturas computacionais de custo proibitivo.

Para a arquitetura de saúde pública no Brasil, em especial no escopo universalista do \textbf{Sistema Único de Saúde (SUS)}, esta democratização representa uma revolução não apenas clínica, mas administrativa e epidemiológica. A redução comprovada de até 58\% no tempo de processamento de exames clínicos, já documentada em projetos-piloto de DT no Brasil~\cite{frota2025sus, abinc2026brasil}, demonstra que a tecnologia possui aderência imediata a cenários de alto fluxo populacional.

Ademais, ao adotar ecossistemas de código aberto, o SUS ganha autonomia soberana sobre os dados (preservando o sigilo previsto na LGPD) e a capacidade de treinar algoritmos focados no genótipo e fenótipo específicos da população latino-americana, mitigando o viés algorítmico importado de sistemas internacionais genéricos.
